% =============================================================================
% Data Engineering
% =============================================================================
\label{chap:data_engineering}

The development of an automated seismic catalog requires robust data 
engineering infrastructure capable of handling heterogeneous data sources, 
legacy file formats, and large-scale data volumes. This section describes the 
data engineering framework developed, encompassing data acquisition, 
multi-format parsing, catalog aggregation, and visualization components.

\section{Software Architecture}

The data engineering pipeline follows a modular, object-oriented design 
implemented in Python, leveraging established seismological libraries and 
modern data science tools. The architecture is organized around a central 
constants module that provides configuration parameters, utility functions, 
and standardized interfaces used throughout the processing chain.

\subsection{Central Configuration Module}
The \texttt{ogsconstants} module serves as the central repository for all 
constants, configuration parameters, and utility functions used throughout the 
pipeline. This design pattern ensures consistency across all processing stages 
and facilitates maintenance. The module provides:

\begin{itemize}
  \item \textbf{Global Configuration}: MPI rank and size for parallel 
  processing, GPU device allocation for machine learning inference, and 
  numerical epsilon values for floating-point comparisons

  \item \textbf{Date/Time Standards}: Standardized format strings following 
  Python's \texttt{strftime} conventions (e.g., \texttt{YYMMDD}, 
  \texttt{YYYYMMDD}, ISO 8601 formats), along with time delta constants for 
  event detection, pick matching, and data segmentation

  \item \textbf{Seismic Phase Identifiers}: Standard labels for P-wave 
  (primary) and S-wave (secondary) arrivals, probability thresholds for 
  machine learning detections, and SEED naming conventions for channel 
  identification

  \item \textbf{Data Column Headers}: Standardized column names for pandas 
  DataFrames, ensuring consistent data structures across all parsing and 
  analysis modules

  \item \textbf{Geographic Definitions}: Polygon boundaries defining the OGS 
  study region in Northeastern Italy, geographic zone codes (Friuli, Veneto, 
  Slovenia, Alto Adige, etc.), and event type classifications

  \item \textbf{FDSN Client Endpoints}: URLs for international data centers 
  (INGV, IRIS, GFZ, ETH, ORFEUS, GEOFON) and OGS-specific internal servers, 
  with a priority-ordered default client list
\end{itemize}

\subsection{Weight Conversion}
The pipeline implements the standard weight conversion table for pick quality 
assessment. Table~\ref{tab:h71_weights} presents the mapping between 
analyst-assigned weight codes and corresponding timing uncertainties in 
seconds.

\begin{table}[h]
  \centering
  \caption{Weight codes and corresponding timing uncertainties used by the OGS}
  \label{tab:h71_weights}
  \begin{tabular}{ccc}
    \hline
    \textbf{Weight} & \textbf{Uncertainty (s)} & \textbf{Interpretation} \\
    \hline
    0 & $<$ 0.01 & Impulsive onset, very clear \\
    1 & 0.01 -- 0.04 & Clear onset \\
    2 & 0.04 -- 0.20 & Fairly clear onset \\
    3 & 0.20 -- 1.00 & Emergent onset \\
    4 & $\geq$ 1.00 & Poor quality pick \\
    \hline
  \end{tabular}
\end{table}

\section{Data Acquisition}

The data acquisition component provides automated retrieval of seismic 
waveforms from distributed FDSN (Federation of Digital Seismograph Networks) 
web services. The implementation supports both rectangular and circular domain 
specifications for geographic selection.

\subsection{FDSN Mass Downloader}
The acquisition module leverages ObsPy's \texttt{MassDownloader} facility to 
efficiently retrieve continuous waveform data from multiple data centers. The 
system supports:

\begin{itemize}
  \item \textbf{Rectangular Domains}: Geographic bounding boxes specified by 
  minimum/maximum longitude and latitude coordinates, suitable for regional 
  studies with well-defined spatial boundaries

  \item \textbf{Circular Domains}: Center point with minimum and maximum radii, 
  useful for studies focused on specific source regions or volcanic systems

  \item \textbf{Multi-Client Querying}: Automatic fallback through a 
  prioritized list of FDSN data centers (OGS internal, INGV, GFZ, IRIS, ETH, 
  ORFEUS) to maximize data availability
\end{itemize}

\subsection{Temporal Segmentation}
To prevent server timeouts and enable efficient parallel downloading, the 
acquisition module segments requests on a daily basis. For each day in the 
requested date range, the system:

\begin{enumerate}
  \item Computes the appropriate time window (full day or centered clip)
  \item Constructs a hierarchical output directory structure organized by 
  year, month, and day
  \item Applies channel priority filters (HH, EH, HN, HG, details in section 
  \ref{subsec:ChannelPriorities}) to select the highest-quality available data.
  \item Enforces minimum interstation distance constraints to avoid 
  redundant co-located recordings
\end{enumerate}

The directory organization follows the pattern:
\begin{verbatim}
{base_directory}/{YYYY}/{MM}/{DD}/
\end{verbatim}
facilitating efficient data access and incremental processing.

\subsection{Channel Selection Priorities}
\label{subsec:ChannelPriorities}
The acquisition system implements a channel priority scheme optimized for 
regional earthquake monitoring:

\begin{enumerate}
  \item \textbf{HH[ZNE]}: High-gain broadband sensors (primary)
  \item \textbf{EH[ZNE]}: Extremely short-period sensors
  \item \textbf{HN[ZNE]}: High-gain accelerometers
  \item \textbf{HG[ZNE]}: High-gain seismometers
\end{enumerate}

Location code priorities ensure consistent station identification across 
networks with varying metadata conventions.

\section{Multi-Format Catalog Parsing}

The OGS seismic monitoring network has accumulated decades of earthquake 
observations stored in standard file formats. The data engineering 
framework provides a unified interface for parsing these data sources through 
an extensible class hierarchy.

\subsection{Abstract Base Class Design}
The \texttt{OGSDataFile} class serves as an abstract base for all 
format-specific parsers. It provides:

\begin{itemize}
  \item \textbf{Regex-Based Extraction}: Configurable regular expression 
  patterns for extracting seismic picks (phase arrivals) and events 
  (earthquake hypocenters) from text-based data files

  \item \textbf{Geographic Filtering}: Polygon-based spatial filtering using 
  matplotlib's Path containment testing to select events within the region 
  of interest

  \item \textbf{Temporal Filtering}: Date range subsetting for extracting 
  specific time periods from multi-year catalog files

  \item \textbf{Parquet Output}: Persistence of parsed data in Apache Parquet 
  columnar format for efficient storage and fast I/O operations

  \item \textbf{Debug Utilities}: Progressive regex testing to identify 
  which capture group fails during parsing, facilitating format adaptation 
  and error diagnosis
\end{itemize}

Subclasses must define format-specific regex patterns 
(\texttt{RECORD\_EXTRACTOR\_LIST} for picks, \texttt{EVENT\_EXTRACTOR\_LIST} 
for events) and implement the \texttt{read()} method for file parsing.

\subsection{Supported File Formats}
The framework supports four primary OGS catalog and input phase formats, each 
handled by a dedicated parser class. Table~\ref{tab:file_formats} summarizes 
the supported formats and their contents.

\begin{table}[h]
\centering
\caption{Supported OGS catalog and input phase formats and their parser
         implementations}
\label{tab:file_formats}
\begin{tabular}{llp{7cm}}
\hline
\textbf{Extension} & \textbf{Parser Class} & \textbf{Content Description} \\
\hline
.hpl & \texttt{DataFileHPL} & Catalog and input phase picks with hypocenter 
  locations, magnitude estimates, quality metrics, and analyst notes \\
.dat & \texttt{DataFileDAT} & Input phase picks (P and S wave arrival times) in 
  fixed-width format with onset quality, polarity, and weight codes \\
.txt & \texttt{DataFileTXT} & Hypocenter location with local magnitude (ML) 
  and duration magnitude (MD) estimates \\
.pun & \texttt{DataFilePUN} & Legacy punch card format containing event 
  hypocenters with quality metrics \\
\hline
\end{tabular}
\end{table}

\subsection{HPL Format Parser}
The HPL (Hypocenter Location) format is the most complete OGS catalog format, 
containing comprehensive event and pick information. The parser extracts:

\begin{itemize}
  \item \textbf{Event Headers}: Origin time, latitude, longitude, depth, 
  magnitude, number of picks, azimuthal gap, RMS residual, horizontal and 
  vertical errors, quality metric

  \item \textbf{Phase Records}: Station code, P-wave onset quality and 
  polarity, P-wave weight, P-wave arrival time, S-wave onset and weight, 
  S-wave arrival time

  \item \textbf{Metadata}: Geographic zone classification (Friuli, Veneto, 
  Slovenia, etc.), event type (local earthquake, explosion, landslide), 
  localization status, analyst notes
\end{itemize}

\subsection{DAT Format Parser}
The DAT format stores phase picks in a fixed-width column structure derived 
from legacy mainframe systems. Each record contains:

\begin{itemize}
  \item Station code (4 characters, right-padded)
  \item P-wave onset quality indicator (e=emergent, i=impulsive, ?=uncertain)
  \item P-wave first motion polarity (+/-/c/d for compression/dilatation)
  \item P-wave weight (0-4 quality scale)
  \item Date-time in compact YYMMDDHHMM format
  \item P-wave arrival time in seconds and centiseconds
  \item Optional S-wave data (time, onset, polarity, weight)
  \item Geographic zone and event type codes
\end{itemize}

\subsection{TXT Format Parser}
The TXT format contains event magnitude information in a simple text structure. 
The parser extracts:

\begin{itemize}
  \item Local magnitude (ML) estimates
  \item Duration magnitude (MD) estimates
  \item Geographic zone and event type codes
\end{itemize}

\subsection{PUN Format Parser}
The PUN (punch card) format is a legacy output format containing event 
summaries. The parser extracts coordinates in degrees-minutes format and 
converts to decimal degrees, along with quality metrics (GAP, DMIN, RMS, 
ERH, ERZ) and quality classification codes.

\subsection{Catalog Aggregation and Merging}

The \texttt{DataCatalog} class provides multi-format support and catalog merging 
capabilities, automatically dispatching to the appropriate format-specific parser 
based on file extension.

\subsection{Automatic Format Detection}
The aggregator maintains a mapping between file extensions and parser classes:

\begin{verbatim}
.hpl → DataFileHPL (hypocenter locations)
.dat → DataFileDAT (phase picks)
.txt → DataFileTXT (magnitude information)
.pun → DataFilePUN (punch card events)
\end{verbatim}

This design allows transparent processing of mixed-format input directories 
without explicit format specification.

\subsection{Merge Logic}
When the merge option is enabled, the aggregator consolidates data from 
multiple input files into a unified catalog:

\begin{enumerate}
  \item \textbf{Picks}: Simple concatenation from all input files, with 
  consistent column naming and data type enforcement

  \item \textbf{Events}: Outer join on time and location, with format-specific 
  attribute handling:
  \begin{itemize}
    \item TXT files contribute magnitude data (ML, MD)
    \item PUN files contribute hypocenter locations
    \item HPL files provide primary event information and associated picks
  \end{itemize}
\end{enumerate}

\subsection{Output Organization}
Merged catalogs are persisted in a date-based directory structure:

\begin{verbatim}
{output}/.all/assignments/{YYYY}-{MM}-{DD}  (picks)
{output}/.all/events/{YYYY}-{MM}-{DD}       (events)
\end{verbatim}

The use of Apache Parquet format provides efficient columnar storage with 
built-in compression, enabling fast analytical queries on large catalogs.



\begin{comment}
\section{Visualization Framework}

Although the visualization utilities described below constitute 
general-purpose tools applicable across multiple stages of the workflow, 
they are presented here for completeness as they were developed as part 
of the data engineering infrastructure.

The visualization module provides a comprehensive set of plotting utilities 
optimized for seismic data analysis and catalog comparison.

\subsection{Base Plotter Class}
All visualization classes inherit from a common base that provides:
\begin{itemize}
  \item Consistent figure sizing and font configuration
  \item Automated output path construction
  \item High-resolution (300 DPI) figure export in multiple formats 
  (PNG, PDF, EPS)
  \item OGS institutional color palette (blue: \#163771, pink: \#E4007C, 
  green: \#00e468, orange: \#FF8C00)
\end{itemize}

\subsection{Specialized Plotters}
The framework includes purpose-built visualization classes:

\begin{itemize}
  \item \textbf{Line Plotter}: Time series visualization with multi-trace 
  overlay support

  \item \textbf{Event Plotter}: Waveform section displays organized by 
  epicentral distance, with automatic trace normalization, bandpass filtering, 
  and pick annotation

  \item \textbf{Stream Plotter}: Individual waveform visualization with 
  channel labeling

  \item \textbf{Day Plotter}: Cumulative event counts over time for 
  catalog completeness assessment

  \item \textbf{Map Plotter}: Geographic visualization using Cartopy 
  projections, with support for:
  \begin{itemize}
    \item Rectangular and circular study region boundaries
    \item OGS polygon region overlay
    \item Station and event scatter plots with custom markers
    \item Inset regional context maps
    \item Coastlines, borders, and ocean features
  \end{itemize}

  \item \textbf{Scatter Plotter}: Generic scatter plot with aspect ratio 
  control and multi-series support

  \item \textbf{Histogram Plotter}: Distribution visualization with 
  optional curve fitting, suitable for magnitude-frequency analysis

  \item \textbf{Confusion Matrix Plotter}: Performance evaluation displays 
  using scikit-learn's \texttt{ConfusionMatrixDisplay}
\end{itemize}

\subsection{Geodetic Distance Calculations}
The visualization module includes optimized vectorized functions for computing 
hypocentral distances:

\begin{equation}
  d = \sqrt{d_{\text{horizontal}}^2 + \Delta z^2}
\end{equation}

where $d_{\text{horizontal}}$ is computed using the WGS84 ellipsoid, enabling 
accurate distance-ordered waveform section displays.
\end{comment}

\section{Data Storage Strategy}

The data engineering framework employs a hierarchical storage strategy 
optimized for both archival and analytical access patterns.

\subsection{Raw Data Storage}
Waveform data are stored in MiniSEED format (the seismological community 
standard) organized by date hierarchy. This structure facilitates:
\begin{itemize}
  \item Incremental acquisition without file locking conflicts
  \item Date-based parallel processing
  \item Efficient archival and backup operations
\end{itemize}

\subsection{Processed Data Storage}
Catalog data (picks and events) are stored in Apache Parquet format, providing:
\begin{itemize}
  \item Columnar compression (typically 5-10x reduction)
  \item Predicate pushdown for efficient filtering
  \item Schema evolution support for format updates
  \item Direct integration with pandas and Spark
\end{itemize}

\subsection{Intermediate Results}
Symbolic links are used to share intermediate results between processing 
stages, reducing storage requirements on shared HPC filesystems while 
maintaining clear provenance tracking.

\section{Quality Assurance}

The data engineering framework incorporates multiple quality assurance 
mechanisms:

\begin{itemize}
  \item \textbf{Input Validation}: Path existence checks, file extension 
  verification, and date range validation before processing

  \item \textbf{Parse Error Logging}: Detailed logging of unparseable 
  records with line numbers and partial match information for debugging

  \item \textbf{Progressive Regex Debugging}: Automated identification 
  of failing capture groups when parsing malformed input

  \item \textbf{Geographic Filtering}: Polygon containment testing to 
  exclude events outside the region of interest

  \item \textbf{Duplicate Detection}: Temporal proximity checking during 
  catalog merging to identify and flag potential duplicates
\end{itemize}

% =============================================================================
% End of Data Engineering Chapter
% =============================================================================